% macros.tex -- notation used throughout the report.
%
% Conventions follow notes/model.md and notes/diagnostics.md. Change a symbol here, not in the
% sections. Symbols that clash with the notes are renamed here on purpose:
%   - the tangent kernel is \TK (Theta), because K is the number of cores (the notes call it K_t);
%   - aligned-init amplitudes are q_{v,r} (\amp), because a_r is the teacher-component overlap;
%   - the gradient-flow preconditioner is diag(eta) (\Pre), because D is the set of differing sites;
%   - the rooted tree is blackboard T, because \mathcal{T}(Z) is the first-escape prefactor;
%   - event widths are omega_r, because w is the weight ratio s_2/s_1.

% ---- general ---------------------------------------------------------------------------------
\newcommand{\R}{\mathbb{R}}
\newcommand{\E}{\mathbb{E}}
\newcommand{\Normal}{\mathcal{N}}
\newcommand{\T}{^{\top}}                          % transpose
\newcommand{\Id}{I}                               % identity matrix
\newcommand{\iid}{i.i.d.\@\xspace}
\DeclareMathOperator{\tr}{tr}
\DeclareMathOperator{\rank}{rank}
\DeclareMathOperator{\diag}{diag}
\DeclareMathOperator{\GL}{GL}
\DeclareMathOperator*{\argmin}{arg\,min}
\DeclarePairedDelimiter{\abs}{\lvert}{\rvert}
\DeclarePairedDelimiter{\norm}{\lVert}{\rVert}
\DeclarePairedDelimiterX{\ip}[2]{\langle}{\rangle}{#1,\,#2}
\newcommand{\normF}[1]{\norm{#1}_{F}}             % Frobenius norm
\newcommand{\normmax}[1]{\norm{#1}_{\max}}        % largest absolute entry
\newcommand{\ipS}[2]{\ip{#1}{#2}_{\Sig}}          % Sigma-weighted inner product
\newcommand{\normS}[1]{\norm{#1}_{\Sig}}          % Sigma-weighted norm
\newcommand{\dd}{\mathrm{d}}                      % differential
\newcommand{\ddt}{\frac{\dd}{\dd t}}

% ---- inputs and data ---------------------------------------------------------------------------
\newcommand{\Ns}{N}                               % number of sites
\newcommand{\pd}[1][]{d_{#1}}                     % physical (local feature) dimension of site i
\newcommand{\xs}[1]{x^{(#1)}}                     % input at site i
\newcommand{\feat}[1]{\phi_{#1}}                  % local feature map of site i
\newcommand{\Feat}{\Phi}                          % product feature tensor
\newcommand{\Sig}{\Sigma}                         % product covariance Sigma_1 (x) ... (x) Sigma_N
\newcommand{\Sigi}[1]{\Sigma_{#1}}                % covariance of site i
\newcommand{\nsamp}{n}                            % training-set size
\newcommand{\batch}{B}                            % minibatch size
\newcommand{\nprobe}{P}                           % size of the kernel probe set

% ---- tree and tensor network ---------------------------------------------------------------------
\newcommand{\tree}{\mathbb{T}}                    % rooted tree (\mathcal{T}(Z) is the first-escape prefactor)
\newcommand{\nodes}{\mathcal{V}}                  % its nodes (= cores)
\newcommand{\edges}{\mathcal{E}}                  % its edges (= bonds)
\newcommand{\Kc}{K}                               % number of cores
\newcommand{\vroot}{\mathrm{root}}                % the root node
\newcommand{\pa}[1]{p(#1)}                        % parent of node v
\newcommand{\ch}[1]{\mathrm{ch}(#1)}              % children of node v
\newcommand{\sites}[1]{S_{#1}}                    % sites attached to node v
\newcommand{\below}[1]{A_{#1}}                    % sites in the subtree below node v
\newcommand{\edge}[1]{e_{#1}}                     % edge between node v and its parent
\newcommand{\G}[1]{G_{#1}}                        % core of node v
\newcommand{\Gm}[2]{G_{#1,(#2)}}                  % core v matricized with its e-leg as rows
\newcommand{\bd}[1]{\chi_{#1}}                    % bond dimension of edge e_v
\newcommand{\fanin}[1]{F_{#1}}                    % fan-in of core v
\newcommand{\params}{\theta}                      % all parameters (cores)
\newcommand{\W}{W}                                % tensor represented by the TTN
\newcommand{\Wt}{W^{*}}                           % teacher tensor
\newcommand{\Wz}{W_{0}}                           % tensor at initialization
\newcommand{\mat}[2]{#1^{\langle #2\rangle}}      % matricization of a tensor across edge e
\newcommand{\Sub}[1]{T_{#1}}                      % contraction of the subtree below v

% ---- initialization ------------------------------------------------------------------------------
\newcommand{\isc}{\sigma}                         % per-core init scale
\newcommand{\iscv}[1]{\sigma_{#1}}                % init scale of core v
\newcommand{\fsc}{\alpha}                         % function-space init scale sigma^K
\newcommand{\amp}[2]{q_{#1,#2}}                   % aligned-init amplitude of core v on component r
\newcommand{\prodamp}[1]{\Pi_{#1}}                % product of amplitudes of component r

% ---- teacher -------------------------------------------------------------------------------------
\newcommand{\ncomp}{R}                            % number of CP components
\newcommand{\tw}[1]{s_{#1}}                       % teacher weight of component r
\newcommand{\fac}[2]{u_{#1}^{(#2)}}               % factor of component r at site i
\newcommand{\U}[1]{U_{#1}}                        % rank-one term of component r (unit norm)
\newcommand{\ov}[1]{a_{#1}}                       % overlap with teacher component r
\newcommand{\Lpart}[1]{L_{#1}}                    % partial-fit loss level

% ---- bond spectra --------------------------------------------------------------------------------
\newcommand{\bs}[2]{s^{#1}_{#2}}                  % k-th singular value across edge e
\newcommand{\bst}[2]{s^{#1,*}_{#2}}               % same for the teacher
\newcommand{\rk}[1]{\operatorname{rk}_{#1}}       % effective bond rank of edge e
\newcommand{\rkH}[1]{\operatorname{rk}^{H}_{#1}}  % entropy rank of edge e

% ---- training ------------------------------------------------------------------------------------
\newcommand{\Loss}{L}                             % population loss
\newcommand{\Lemp}{\hat{L}_{\nsamp}}              % empirical loss
\newcommand{\lr}{\eta}                            % learning rate (step size) of GD/SGD
\newcommand{\lrm}[1]{\eta_{#1}}                   % per-core learning-rate multiplier
\newcommand{\mask}{\mathfrak{m}}                  % 0/1 gradient mask (m is |V_D|)
\newcommand{\Pre}{\operatorname{diag}(\boldsymbol{\eta})} % diagonal preconditioner (per-entry rates and mask); D is the site set
\newcommand{\tmax}{t_{\max}}

% ---- symmetries ------------------------------------------------------------------------------
\newcommand{\Q}[1]{Q_{#1}}                        % conserved quantity of bond e

% ---- diagnostics ---------------------------------------------------------------------------------
\newcommand{\TK}{\Theta}                          % tangent kernel (tensor kernel; restricted to the probe set where stated)
\newcommand{\dTK}{\delta_{\TK}}                   % relative kernel change
\newcommand{\dth}{\delta_{\params}}               % relative parameter change
\newcommand{\KTA}{\mathcal{A}}                    % kernel-target alignment
\newcommand{\Jac}{J}                              % Jacobian dW/dtheta
\newcommand{\Hess}{H}                             % Hessian of the loss
\newcommand{\Hperp}{H_{\perp}}                    % gauge-projected Hessian
\newcommand{\sep}[1]{\rho_{#1}}                   % separation ratio of consecutive events
\newcommand{\tev}[1]{t_{#1}}                      % event time
\newcommand{\wev}[1]{\omega_{#1}}                 % event width (w is the weight ratio)
\newcommand{\tmid}[1]{t^{\mathrm{mid}}_{#1}}      % plateau midpoint
